\documentclass[11pt]{article}
\usepackage[margin=1in]{geometry}
\usepackage{natbib}
\usepackage{booktabs}
\usepackage{amsmath}
\usepackage{graphicx}
\usepackage{hyperref}
\usepackage{authblk}

\title{A SUMO E3 Ligase Promotes Long Non-Coding RNA Transcription to Regulate Small RNA-Directed DNA Elimination}

\author[1]{Salman Shehzada}
\author[1]{Kazufumi Mochizuki}
\affil[1]{Institute of Human Genetics, University of Montpellier, CNRS, Montpellier, France}

\date{}

\begin{document}
\maketitle

\begin{abstract}
Small RNAs target their complementary chromatin regions for gene silencing through nascent long non-coding RNAs (lncRNAs). In the ciliated protozoan \textit{Tetrahymena}, the interaction between Piwi-associated small RNAs (scnRNAs) and the nascent lncRNA transcripts from the somatic genome has been proposed to induce target-directed small RNA degradation (TDSD), and scnRNAs not targeted for TDSD later target the germline-limited sequences for programmed DNA elimination. In this study, we show that the SUMO E3 ligase Ema2 is required for the accumulation of lncRNAs from the somatic genome and thus for TDSD and completing DNA elimination to make viable sexual progeny. Ema2 interacts with the SUMO E2 conjugating enzyme Ubc9 and enhances SUMOylation of the transcription regulator Spt6. We further show that Ema2 promotes the association of Spt6 and RNA polymerase II with chromatin. These results suggest that Ema2-directed SUMOylation actively promotes lncRNA transcription, which is a prerequisite for communication between the genome and small RNAs.
\end{abstract}

\section{Introduction}

Small RNAs of approximately 20--30 nucleotides that are complexed with Argonaute family proteins target either mRNAs for post-transcriptional silencing or chromatin regions for transcriptional and co-transcriptional gene silencing \citep{volpe2002,verdel2004,sienski2012}. For the latter process, small RNAs are generally considered to recognize their genomic targets via nascent long non-coding RNAs (lncRNAs) to induce heterochromatin formation. Therefore, chromatin regions that are targeted for silencing by small RNAs must be paradoxically transcribed to provide nascent lncRNAs.

In fission yeast, small interfering RNAs (siRNAs) mediate the deposition of histone H3 lysine 9 methylation (H3K9me) at pericentromeric repeats \citep{volpe2002,verdel2004}. While the HP1 protein Swi6 binds to H3K9me2/3 for transcriptional silencing, phosphorylation of histone H3 serine 10 at the M phase of the cell cycle evicts Swi6, passively allowing lncRNA transcription \citep{zhang2008,kloc2008}. In contrast, the HP1 paralog Rhino in fruit flies, which is specifically enriched at PIWI-interacting RNA (piRNA) clusters by binding to H3K9me2/3 \citep{lethomas2014}, actively recruits a dedicated transcription machinery to produce piRNA precursor lncRNAs \citep{klattenhoff2009,andersen2017}. SUMOylation of the transcriptional co-repressor Panoramix on piRNA-targeted chromatin connects the piRNA pathway to the cellular heterochromatin machinery \citep{batki2022}. In plants, the RNA-directed DNA methylation (RdDM) pathway also employs specialized transcription factors to produce scaffold lncRNAs \citep{law2011,law2013}.

In the ciliated protozoan \textit{Tetrahymena thermophila}, programmed DNA elimination removes internal eliminated sequences (IESs) from the developing somatic macronucleus (MAC) during sexual reproduction (conjugation). This process requires scan RNAs (scnRNAs), a class of Piwi-associated small RNAs produced from the germline micronucleus (MIC) during early conjugation stages \citep{mochizuki2002analysis}. The interaction between scnRNAs and nascent lncRNA transcripts from the parental MAC (pMAC-lncRNAs) has been proposed to induce target-directed small RNA degradation (TDSD), selectively retaining only IES-complementary scnRNAs that later direct DNA elimination \citep{aronica2008study,noto2018}. The transcription of pMAC-lncRNAs is bidirectional and produces long double-stranded RNAs (dsRNAs) \citep{chalker2001,woo2016,mochizuki2005,malone2005}. A meiosis-specific Spt5 homolog has been implicated in non-coding transcription in this system \citep{kawakami2017}. However, the molecular mechanisms that actively promote lncRNA transcription in the parental MAC have remained unclear.

Here, we report that the SUMO E3 ligase Ema2 is essential for the accumulation of pMAC-lncRNAs and, consequently, for TDSD and completion of DNA elimination. We demonstrate that Ema2 interacts with the SUMO E2 conjugating enzyme Ubc9, enhances SUMOylation of the transcription elongation factor Spt6, and promotes the chromatin association of both Spt6 and RNA polymerase II (RNAPII). Unexpectedly, a SUMOylation-defective Spt6 mutant retains the ability to support lncRNA transcription, suggesting that Ema2-directed SUMOylation may promote lncRNA transcription through additional targets or mechanisms beyond Spt6 SUMOylation alone.

\section{Results}

\subsection{Ema2 is exclusively expressed during conjugation and localized in the MAC}

As part of a systematic investigation into genes highly upregulated during conjugation \citep{loidl2021}, we explored the function of \textit{EMA2} (TTHERM\_00113330). \textit{EMA2} mRNA is exclusively expressed during conjugation (Figure~1A). The encoded Ema2 protein tagged with HA at the endogenous \textit{EMA2} locus was not detectable in vegetative cells and first appeared in the MAC during conjugation (Figure~1B, 3 hours post-induction of mating [hpm] and 6 hpm). At the onset of new MAC development, Ema2 disappeared from the parental MAC and appeared in the new MAC (Figure~1B, 8 hpm), which later faded away (12 and 14 hpm). The conjugation-specific expression and the localization switch from the parental to the new MAC are reminiscent of factors involved in DNA elimination such as the Piwi protein Twi1 \citep{mochizuki2002analysis,liu2007h3k9,noto2010}.

\subsection{Ema2 is required for completing DNA elimination}

DNA elimination of exconjugants at 36 hpm was analyzed by DNA fluorescent \textit{in situ} hybridization (FISH) using probes complementary to the transposable element Tlr1 \citep{wuitschick2002}. DNA elimination is completed by approximately 14--18 hpm in wild-type cells \citep{austerberry1984,mutazono2019}, and the Tlr1 element was detected only in the MICs of exconjugants from wild-type cells (Figure~2A). In contrast, the Tlr1 element was detected in both the MICs and the MACs of exconjugants from the \textit{EMA2} somatic KO strains \citep{shehzada2022} (Figure~2A). The intensity of the FISH signal in the new MACs was lower in the \textit{EMA2} KO cells than in \textit{TWI1} KO cells (Figure~2B), in which DNA elimination is known to be completely blocked \citep{noto2015}. Therefore, DNA elimination was partially blocked in the exconjugants of \textit{EMA2} KO cells. Consistent with the requirement of DNA elimination for the viability of sexual progeny \citep{cheng2010,vogt2013}, \textit{EMA2} KO cells did not produce viable progeny (Figure~2C).

\subsection{Ema2 is required for target-directed small RNA degradation}

We next asked whether \textit{EMA2} is involved in TDSD. The production of scnRNAs occurs from IESs and their surrounding macronucleus-destined sequences (MDS) regions in the MIC at the early conjugation stages, and scnRNAs complementary to the MAC genome (i.e., MDSs) are subjected to TDSD in the parental MAC in the mid-conjugation stages, resulting in the selective retention of IES-complementary scnRNAs \citep{schoeberl2012}. Using the Mi-9 probe complementary to a repetitive sequence found in both IESs and MDSs \citep{aronica2008study}, northern blot analysis showed that in wild-type cells, Mi-9-complementary scnRNAs were detected at 3 hpm, reduced at 4.5 hpm, and became undetectable at 6 hpm due to TDSD (Figure~3A). In contrast, these scnRNAs remained detectable at 6 hpm in \textit{EMA2} KO cells, suggesting that Ema2 is required for TDSD.

Small RNA sequencing at 3 and 8 hpm confirmed this result at the genome-wide level. Although scnRNAs that mapped to many MDS tiles were greatly reduced from 3 to 8 hpm due to TDSD in wild-type cells (Figure~3B), those in \textit{EMA2} KO cells were only slightly reduced by 8 hpm. This TDSD defect in \textit{EMA2} KO cells was milder than that in \textit{EMA1} KO cells, in which the amount of MDS-complementary scnRNAs remained constant until 8 hpm (Figure~3B), indicating that Ema1 and Ema2 act differently in TDSD. In contrast, the amounts of scnRNAs complementary to IES tiles at 3 and 8 hpm were comparable in all strains, consistent with previous observations that only MDS-complementary scnRNAs are targeted for TDSD.

\subsection{Ema2 acts as a SUMO E3 ligase}

Ema2 contains an atypical SP-RING domain (Figure~4A) in which the first cysteine of the zinc ion-binding residues in typical counterparts is replaced by histidine, and some stabilizer residues \citep{duan2009,yunus2009} are likely missing. SUMOylation is catalyzed by the sequential actions of E1, E2, and in most cases E3 enzymes, and thus Ema2 should interact with the E2 enzyme if it acts in SUMOylation. We found that Ema2 can directly interact with the \textit{Tetrahymena} SUMO E2 enzyme Ubc9 \textit{in vitro} (Figure~4B), confirming that Ema2 is a bona fide SUMO E3 ligase.

To examine the role of Ema2 in SUMOylation, we expressed HA-tagged Smt3 (HA-Smt3) in an \textit{EMA2} KO strain. SUMO is solely encoded by \textit{SMT3} in \textit{Tetrahymena} \citep{nasir2015}, and HA-Smt3 could replace the essential function of endogenous Smt3 (Figure~EV1). Total proteins harvested at 4.5 and 6 hpm showed that high molecular weight SUMOylated proteins (mainly $>$200 kDa) were detected in the WT cross, and they were reduced to approximately 50\% in the \textit{EMA2} KO cross (Figure~4C). These results indicate that Ema2 is the major SUMO E3 ligase during the mid-conjugation stages. The requirement of protein SUMOylation in DNA elimination was previously demonstrated in \textit{Paramecium} \citep{matsuda2006}, suggesting that involvement of a SUMO pathway in DNA elimination is conserved among ciliates.

\subsection{Ema2 promotes SUMOylation of Spt6}

To identify the SUMOylation target(s) of Ema2, we introduced a construct expressing His-tagged Smt3 and concentrated SUMOylated proteins using nickel-NTA beads for mass spectrometry analysis. Although most proteins were detected similarly between the WT and \textit{EMA2} KO crosses, Spt6, the most abundantly detected protein in the WT cross, was detected at a much lower level in the \textit{EMA2} KO cross (Figure~5A), suggesting that Spt6 is the major Ema2 target for SUMOylation.

Confirmation by western blotting using HA-tagged Spt6 showed that a slower migrating population of Spt6-HA was detected in the WT cross at both 4.5 and 6 hpm (Figure~5B), while this population was barely detectable in the \textit{EMA2} KO cross. Immunoprecipitation of Spt6-HA followed by western blotting with an anti-Smt3 antibody confirmed that the slower migrating species were SUMOylated Spt6 and that Ema2 is required for the majority of Spt6 SUMOylation during mid-conjugation (Figure~5C).

\subsection{Ema2 is required for the accumulation of lncRNAs in the parental MAC}

Spt6 is a histone chaperone and transcription elongation factor. Because TDSD was proposed to be triggered by the base-pairing interaction between scnRNAs and nascent pMAC-lncRNAs \citep{aronica2008study,noto2018}, we hypothesized that Ema2-dependent Spt6 SUMOylation promotes pMAC-lncRNA transcription. RT-PCR analysis showed that pMAC-lncRNAs were detected in wild-type cells but not in \textit{EMA2} KO cells at 6 hpm for three tested loci, although the control \textit{RPL21} mRNA was detected in both conditions (Figure~6A). The lack of detection was not due to loss of the primer binding sites by alternative DNA elimination (Figure~EV3).

Immunostaining using the long dsRNA-specific J2 antibody \citep{schoenborn1991} detected long dsRNAs in the MIC at 3 hpm, in the parental MAC at 6 hpm, and in the new MAC at 9 hpm in wild-type cells (Figure~6B). In \textit{EMA2} KO cells, long dsRNAs were detected in the MIC at 3 hpm and in the new MAC at 9 hpm but were undetectable in the parental MAC at 6 hpm, indicating that Ema2 is specifically required for the accumulation of pMAC-lncRNAs. In contrast, EMA1 was dispensable for the accumulation of pMAC-lncRNAs.

We also found that H3K27me3 was undetectable in the parental MAC in the \textit{EMA2} KO cells (Figure~6C), consistent with the loss of PRC2 recruitment by the scnRNA-lncRNA interaction.

\subsection{Ema2 facilitates chromatin association of Spt6 and RNA polymerase II}

Immunofluorescence staining showed that Spt6-HA and Rpb3 (the third largest subunit of RNAPII) were similarly detected in the MAC in both WT and \textit{EMA2} KO crosses at 4.5 hpm (Figure~7A). However, cell fractionation experiments revealed that while Spt6-HA and Rpb3 were detected in both soluble and chromatin-bound fractions in the WT cross, their levels in the chromatin-bound fraction were greatly reduced in the \textit{EMA2} KO cross (Figure~7B). Importantly, SUMOylated Spt6-HA was detected only in the chromatin-bound fraction. The chromatin association of Spt6 and RNAPII does not require RNA, as they were retained in the insoluble fraction after RNase A treatment (Figure~7C). These findings were validated in replicated experiments (Figures~EV5, EV6), demonstrating that Ema2 promotes the chromatin association of both Spt6 and RNAPII in the parental MAC.

\subsection{Spt6 SUMOylation is dispensable for pMAC-lncRNA transcription}

To directly examine the importance of Spt6 SUMOylation, we produced Spt6 mutants with lysine-to-arginine substitutions. Among the mutants tested, HA-Spt6-C-KR, which has substitutions in the C-terminal region, showed no detectable SUMOylation (Figure~8B), representing a SUMOylation-defective mutant. Surprisingly, cell fractionation experiments showed that HA-Spt6 and RNAPII were detected in the chromatin-bound fraction in \textit{HA-SPT6-C-KR} cells similar to \textit{HA-SPT6-WT} cells (Figure~8C), and long dsRNAs were detected in the parental MAC of the C-KR mutant (Figure~8D). These results suggest that Spt6 SUMOylation per se is not required for lncRNA transcription in the parental MAC. Nonetheless, exconjugants from the \textit{HA-SPT6-C-KR} cells showed a mild DNA elimination defect (Figure~8E), indicating that Ema2-directed Spt6 SUMOylation plays a role in DNA elimination other than promoting pMAC-lncRNA transcription.

\section{Discussion}

In this study, we demonstrate that the SUMO E3 ligase Ema2 is essential for the accumulation of lncRNAs from the parental MAC in \textit{Tetrahymena}, and consequently for TDSD and programmed DNA elimination. Ema2 directly interacts with the SUMO E2 enzyme Ubc9 and promotes SUMOylation of the transcription elongation factor Spt6. Loss of Ema2 results in reduced chromatin association of both Spt6 and RNAPII in the parental MAC.

The finding that a SUMOylation-defective Spt6 mutant retains the ability to support pMAC-lncRNA transcription suggests that Ema2 may promote lncRNA transcription through additional, yet-to-be-identified SUMOylation targets. The remaining approximately 50\% SUMOylation activity in the absence of Ema2 is likely mediated by other SUMO E3 ligases and/or E3-independent SUMOylation \citep{sampson2001}. Nevertheless, the mild DNA elimination defect observed in the SUMOylation-defective Spt6 mutant indicates that Spt6 SUMOylation contributes to DNA elimination through mechanisms distinct from promoting lncRNA transcription.

The involvement of SUMOylation in small RNA-directed chromatin regulation appears to be a recurring theme across eukaryotes. In \textit{Drosophila}, SUMOylation of the transcriptional co-repressor Panoramix connects the piRNA pathway to heterochromatin machinery \citep{batki2022}. In \textit{C. elegans}, HDAC1 SUMOylation promotes Argonaute-directed transcriptional silencing \citep{hdac2021}. Our findings in \textit{Tetrahymena} reveal that SUMOylation plays an active role in promoting lncRNA transcription, which is the prerequisite for the genome-small RNA communication in the context of programmed DNA elimination.

The SUMOylation pathway's involvement in DNA elimination is conserved among ciliates, as demonstrated previously in \textit{Paramecium} \citep{matsuda2006}. Given the widespread importance of lncRNA transcription in small RNA-directed chromatin regulation across eukaryotes, our identification of Ema2 as a key factor connecting SUMOylation to lncRNA transcription provides fundamental insights into the molecular mechanisms underlying genome-small RNA crosstalk.

\section{Materials and Methods}

\subsection{Strains and cell culture conditions}
The \textit{Tetrahymena thermophila} wild-type strains B2086, CU427, and CU428 and the MIC-defective ``star'' strains B*VI and B*VII were obtained from the Tetrahymena stock center. The production of the \textit{EMA2} KO strains was described previously \citep{shehzada2022}. Cells were grown in SPP medium \citep{gorovsky1975} containing 2\% proteose peptone at 30$^\circ$C. For conjugation, growing cells ($\sim$5--7 $\times$ 10$^5$/mL) of two different mating types were washed, pre-starved ($\sim$12--24 hr), and mixed in 10 mM Tris (pH 7.5) at 30$^\circ$C.

\subsection{Antibodies}
The following antibodies were used: anti-H3K27me3 (Millipore, 07-449); anti-HA HA11 (Covance, clone 6B12); anti-H3K9me3 (Active Motif, 39162); anti-GST (BD Biosciences, 554805); anti-His (Proteintech, 66005-1-Ig); anti-long dsRNA J2 (Jena Bioscience RNT-SCI-10010200); anti-$\alpha$-tubulin 12G10 (Developmental Studies Hybridoma Bank). The anti-Rpb3 antibody was described previously \citep{kataoka2017}. The anti-Smt3 antibody was a gift from Dr. James Forney and described previously \citep{nasir2015}.

\subsection{DNA-FISH}
DNA elimination was analyzed by DNA fluorescent \textit{in situ} hybridization using probes complementary to the Tlr1 transposable element as described previously \citep{wuitschick2002}.

\subsection{Small RNA analysis}
Small RNA sequencing was performed at 3 and 8 hpm. Sequenced 26- to 32-nt small RNAs corresponding to scnRNAs \citep{schoeberl2012} were mapped to genomic tiles of MDSs and IESs.

\subsection{Cell fractionation}
Conjugating cells were incubated with lysis buffer to release cytoplasmic and nucleoplasmic proteins into the soluble fraction (S1). The insoluble fraction was sonicated to release chromatin-bound proteins (S2). For RNase treatment experiments, the insoluble fraction was incubated in fresh lysis buffer with or without RNase A prior to sonication.

\subsection{Mass spectrometry}
SUMOylated proteins were concentrated using nickel-NTA beads from cells expressing His-tagged Smt3 and identified by mass spectrometry. Label-free quantification (LFQ) was used for comparison between WT and \textit{EMA2} KO crosses.

\subsection{GST pull-down assay}
GST alone and GST-tagged Ema2 were recombinantly expressed in \textit{E. coli}, immobilized on glutathione beads, and incubated with His-tagged Ubc9. Retained proteins were analyzed by western blotting.

\bibliographystyle{plainnat}
\bibliography{references}

\end{document}
